\chapter{Conclusion and Future Work}
\label{chap:conclusion}

\section{Thesis Summary}

This thesis developed and physically evaluated \QFlow, a fixed-point FPGA
accelerator for adaptive key-resource-aware route-candidate evaluation in
trusted-relay QKD networks. The system accepts externally generated candidate
descriptors, applies hard eligibility, selects a fixed or adaptive scoring profile,
evaluates candidates with deterministic arithmetic, and returns a selected
identifier and status. Candidate generation, topology management, KMS operation,
SDN orchestration, security validation, and route installation remain control-plane
functions.

The adaptive controller quantizes four features into a 256-state address and selects
one of four scoring profiles. Tabular Q-learning is performed offline. The deployed
reinforcement-learned configuration uses a 256-entry two-bit policy ROM and
minimum-dwell control, and performs deterministic inference without runtime Q-value
updates.

The executed features are minimum key occupancy, bounded bottleneck route quality,
offered request load, and utilization imbalance. The second feature is limited to
its role as an abstract experimental input and is not a stored-classical-key
coherence quantity. This experimental contract is kept separate from the final
network model based on occupancy, generation, consumption, status, utilization,
freshness policy, and external health indicators.

\section{Main Findings}

Five independent hardware evaluation configurations were implemented on a Digilent
Nexys 3 Spartan-6 FPGA under a common interface, clock, replay, and measurement
contract. All configurations exceeded the 100~MHz post-route target. The physical
campaign produced 420 records and 3,780 checked fields, all of which matched the
fixed-point reference exactly.

Configuration H4 achieved a post-route maximum frequency of 102.020~MHz and a mean
normal-case accepted-request-to-done kernel latency of 2.048~$\mu$s. Relative to
configuration H3, configuration H4 changed mean latency by +0.986\% and maximum
frequency by $-0.612$\%, mapped to 7.081\% fewer LUTs, 5.378\% more registers, and
4.961\% more occupied slices. These values describe one implemented mapping and do
not imply a universal resource property of reinforcement learning.

The complete configuration H4 controller, including the learned policy ROM and
minimum-dwell mechanism, produced 24 profile transitions compared with 44 for
configuration H3, a 45.455\% reduction on the common replay. The corresponding
route-quality comparisons were not statistically significant: the configuration H4
versus H2 and H4 versus H3 sign-test $p$-values were 0.2188 and 0.4531, and both
confidence intervals crossed zero. The evidence therefore supports physical
deployability, exactness, and distinct adaptive behavior, but does not establish
route-quality superiority.

The supporting VLSI study showed that selected kernels can pass through generic
synthesis and separate OpenROAD flows. Fixed-coefficient specialization reduced one
scoring kernel from 5,343 to 503 generic cells, while isolated routed areas were
6,451~$\mu\mathrm{m}^{2}$ for SKAG-W1, 22,186~$\mu\mathrm{m}^{2}$ for the
LUT/interpolation arithmetic kernel (repository identifier FDPE-V3), and
3,368~$\mu\mathrm{m}^{2}$ for Pareto-C0. These results indicate a possible future
ASIC path; they are not an integrated ASIC implementation.

\section{Contributions}

The work contributes:

\begin{enumerate}
  \item a classical key-resource model for candidate evaluation that represents
  pool occupancy, generation, consumption, status, utilization, health indicators,
  and policy-controlled freshness;
  \item a bounded fixed-point evaluator with separate hard eligibility and soft
  multi-objective preference;
  \item fixed, rule-adaptive, and offline-learned profile-selection mechanisms;
  \item a deterministic policy-ROM deployment architecture with minimum-dwell
  control;
  \item five independent FPGA configurations and a common physical measurement
  contract; and
  \item reproducible exactness, timing, latency, resource, adaptive-behavior, and
  neutral route-quality evidence.
\end{enumerate}

The principal contribution is thus a reproducible hardware-systems result rather
than a claim of universal routing superiority.

\section{Limitations}

The candidate evaluator was tested on a bounded record format, one FPGA board and
device family, and 84 records per configuration. The replay establishes exact
implementation behavior for those inputs but does not establish network-level
superiority on unseen topologies or live QKD deployments. The policy was trained in
a controlled research environment, and replay rows may contain dependence that
reduces the effective statistical sample size.

No exactly matched CPU/FPGA timing benchmark, activity-calibrated FPGA power
measurement, or CPU/FPGA energy comparison was completed. No live KMS, SDN
controller, or QKD device was integrated, and the host/FPGA link was an experimental
transport rather than a production authenticated interface. The VLSI results concern
separate kernels and do not provide integrated area, full-chip power, foundry
signoff, tapeout, or fabricated performance.

\section{Future Work}

The highest-priority next step is an exactly matched CPU/FPGA benchmark using the
same candidate descriptors, fixed-point contract, request protocol, and timing
boundary. This would support a defensible kernel speed comparison and expose the
effect of software scheduling and data transfer.

Second, FPGA power and energy should be measured with activity-calibrated switching
data and a defined operating point for every configuration. Energy per accepted
candidate-evaluation transaction should be separated from board idle power and host
transport.

Third, the policy should be tested on larger independent trace sets and unseen
topologies. The evaluation should include stronger RL and non-RL baselines, bursts,
key-pool depletion and recovery, failures, changing key-generation rates, and
statistically powered route-quality analysis.

Fourth, a policy-by-dwell ablation should independently vary the policy and
minimum-dwell mechanism. This would determine how much of the observed switching
reduction arises from the learned mapping, the dwell rule, and their interaction.

Fifth, \QFlow\ should be integrated with a live KMS/SDN testbed. The integration must
measure metadata acquisition, candidate generation, host/FPGA communication,
security-policy checks, fallback, and route installation separately from kernel
latency.

Sixth, the host/FPGA interface should be authenticated and hardened. Input
validation, replay protection, policy identity, signed bitstreams where supported,
fault containment, audit logging, and software fallback are important for a
security-sensitive controller.

Finally, an integrated full-core ASIC study can be considered after the mathematical
and interface contracts stabilize. Such work would require complete logic and
physical verification, clock and power design, timing closure, DRC/LVS, packaging,
test planning, and a measured comparison with the FPGA implementation.

\section{Closing Statement}

The experiments show that an offline-learned key-resource-aware controller can be
deployed bit-exactly on a resource-constrained FPGA with deterministic
microsecond-scale candidate-evaluation latency and controlled implementation
overhead. Equally important, the neutral route-quality result defines where the
evidence stops. This combination of physical exactness, transparent trade-offs, and
bounded interpretation provides a sound basis for larger network experiments and
future deployment studies.
